\cotefile{158}
\scanpage{53}
\mspage{7}
% @ top, line 1 | text
Comme corollaire, on a ceci

% @ upper half, lines 2–8 | text
\uline{Proposition 2} Soit $X$ une catégorie \uline{finie} pseudo-cofiltrante, \deleted{\unsure{Donc que|Donc qui}} et soit $X_0$ la \inserted{\unsure{ss-|ss|est}}catégorie ouverte de $X$ formée des \inserted{plus petits} objets de $X$ (NB les hyp. impliquent $X_0\ne\varnothing$), car $X_0$ est $0$-connexe). Supposons que $X_0$ soit $1$-connexe.

% @ middle, lines 9–13 | text
Alors $X$ est cofiltrante (donc tout objet $0$-connexe de $X$ est $\infty$-connexe, et plus \unsure{génér.lt|générlt} asphérique pour tt \deleted{\unsure{dériseur|dériveur}} localiseur fondamental $W$ satisfaisant $W(\unsure{8|\delta})$).

% @ lower third, lines 14–20 and foot | text
\inserted{A fortiori :} \uline{Corollaire} Soit $X$ une catégorie \uline{finie} telle que tout objet $0$-connexe de $X^\wedge$ soit $1$-connexe (p. ex. telle que $X^\circ$ soit tot. $W$-asphérique, où $W$ est un loc. fondamental tel que $W\subset\mathbb{W}_1$). Alors $X$ est cofiltrante, i.e. $X^\circ$ est filtrante (et tt objet de $X^\wedge$ est $\infty$-connexe).
% archive sheet number: 52
